\documentclass[11pt]{article}
\usepackage[letterpaper,top=0.96in,bottom=0.97in,left=1.16in,right=1.16in]{geometry}
\usepackage{lmodern}
\usepackage{amsmath,amssymb,amsthm}
\usepackage[colorlinks=true,urlcolor=blue,linkcolor=black,citecolor=black]{hyperref}
\hypersetup{pdftitle={Counting and stability in graph reconstruction},pdfauthor={Mira}}
\newtheorem{theorem}{Theorem}
\newtheorem{corollary}[theorem]{Corollary}
\newtheorem{lemma}[theorem]{Lemma}
\newtheorem{proposition}[theorem]{Proposition}
\theoremstyle{definition}\newtheorem{definition}[theorem]{Definition}\theoremstyle{plain}
\newcommand{\Aut}{\operatorname{Aut}}
\newcommand{\LWPP}{\mathsf{LWPP}}
\newcommand{\Leg}{\textnormal{\textsc{Legitimate Deck}}}
\setlength{\parindent}{1.2em}
\title{Counting and stability in graph reconstruction}
\author{Mira}
\date{November 2026}
\begin{document}
\maketitle
\begin{abstract}
We give three bounds for the ordinary deck of a finite simple graph, the multiset of its unlabeled one-vertex deletions. First, an $n$-card deck, $n\ge3$, has at most $(n!)^{2/(n-2)}\sim n^2/e^2$ nonisomorphic parents. The proof applies a projection inequality to the edge bits of all labeled parents. This answers the polynomial preimage question of Hemaspaandra, Hemaspaandra, Spakowski and Watanabe and, by their theorem, places recognition of legitimate decks in the counting class $\LWPP$. Second, for all sufficiently large $n$, every $n$-vertex graph is within $O(\log n)$ edge edits of a graph uniquely determined by its deck. Third, for independent edge probabilities in $[\beta,1-\beta]$ with $\beta n/\log n\to\infty$, with probability at least $1-\exp(-c\beta n^2)$ every graph within $c\beta n^2$ edits of the sample is uniquely determined by its deck, for an absolute $c>0$. These two distance bounds follow from a structural constraint: a graph not determined by its deck must have linearly many disjoint pairs of vertices whose neighborhoods differ at only $O(\log n)$ other vertices.
\end{abstract}

\section{The size and location of the ambiguity}
The \emph{vertex deck} of a finite simple undirected graph $G$ is the multiset of isomorphism types of the graphs $G-v$ for $v\in V(G)$. A graph with deck $D$ is called a \emph{parent} of $D$; the deck is \emph{legitimate} if it has a parent. Write $p(D)$ for the number of nonisomorphic parents. The Reconstruction Conjecture asks whether $p(D)=1$ for every legitimate deck of order at least three.

Two quantitative questions survive even if we cannot decide uniqueness. How many different graphs could one deck describe? And where in the space of graphs could a failure of uniqueness occur? They call for different uses of the missing labels. For the first, it suffices to count complete labeled graphs from the supports of their deletions. For the second, we need a structural reason why the deck cannot distinguish a graph: many pairs of vertices must have nearly the same adjacency to all the others.

We measure the distance between two graphs on the same vertex set by the number of unordered pairs on which their edge bits differ. The two metric assertions point in different directions. From an \emph{arbitrary} graph a small, chosen change suffices to reach a reconstructible graph. Around a \emph{random} graph, with high probability even an adversary who has seen the sample cannot reach a nonreconstructible graph without many changes. Neither statement supplies a nonreconstructible graph. Asilis, Chen, Hansen and Teng~\cite{acht} put the first distance question explicitly in their outlook and studied the second in semi-random models.

\subsection*{Counting the possible parents}
Hemaspaandra, Hemaspaandra, Spakowski and Watanabe~\cite[Conj.~4.5]{hhsw} asked about bounding $p(D)$ by a function $q$ of the order and singled out the case of polynomial $q$. A single card on $n-1$ vertices gives the immediate bound $2^{n-1}$: each parent can be constructed by choosing one neighborhood on the card. For a bounded minimum degree $k$, one can choose a corresponding card and enumerate $O(n^k)$ extensions~\cite[proof of Thm.~6.1]{kh}. Counting only neighborhoods of the prescribed degree, and taking account of the card's symmetries, improves this estimate for individual decks; compare Wang~\cite[Thms.~14.4 and 15.2]{wang}. We instead use the $n$ deletions simultaneously.

\begin{theorem}\label{thm:main}
For every legitimate deck $D$ of order $n\ge3$,
\[
p(D)\le (n!)^{2/(n-2)}=(1+o(1))\,n^2/e^2 .
\]
More precisely, if $b(D)$ is the number of graphs on $[n-1]$ whose unlabeled type is a card of $D$, then $p(D)\le b(D)^{2/(n-2)}$.
\end{theorem}

The distinction between $p(D)$ and $b(D)$ is useful. The former counts \emph{unlabeled parents}; the latter counts \emph{labeled cards}, ignoring repetitions of a card type. Given a single parent, every one of those $b(D)$ cards can be seen at any designated hole simply by relabeling the parent. Thus a large card family forces \emph{each} parent to have many labeled representatives. The proof compares this lower bound with the number of full graphs the card family can support.

\section{Proof of the bound}
We recall the finite form of Shearer's inequality~\cite{shearer}; compare the uniform-cover method for hereditary hypergraph families of Bollob\'as and Thomason~\cite{bt}. For a set of words $U$ on finitely many coordinates, write $P_A(U)$ for the set of restrictions of its words to a set $A$ of coordinates.
\begin{lemma}\label{lem:proj}
If $U$ is finite and nonempty, and each coordinate belongs to at least $r$ of the sets $A_1,\ldots,A_m$, then
\[
 |U|^r\le\prod_{i=1}^{m}|P_{A_i}(U)| ,
\]
where the product counts each projection separately.
\end{lemma}
\begin{proof}
Let $X$ be uniform on $U$. Order the coordinates and write $X_{<j}$ for those preceding $j$. In the chain rule for $H(X_{A_i})$, the term for $j\in A_i$ is conditioned only on the earlier coordinates in $A_i$. Since conditioning reduces entropy,
\[
\sum_i H(X_{A_i})\ \ge\ \sum_i\sum_{j\in A_i}H(X_j\mid X_{<j})
\ \ge\ r H(X).
\]
Use $H(X)=\log_2|U|$ and $H(X_{A_i})\le\log_2|P_{A_i}(U)|$, and exponentiate.
\end{proof}

\begin{proof}[Proof of Theorem~\ref{thm:main}]
Work on the common vertex set $[n]$. For each parent isomorphism class $G$ let $O_G$ be its set of labeled representatives. Put $U=\bigcup_G O_G$; the union is disjoint. Write $\pi_i$ for restriction to $[n]\setminus\{i\}$, that is, deletion of the label $i$ without changing the surviving labels. For each $G$ and $i$, $\pi_i(O_G)$ is exactly the set of all labeled copies, on this fixed surviving set, of the card types in $D$. Indeed, for any one such copy, choose a vertex of $G$ whose deletion has that type, label the remaining vertices by an isomorphism to the copy, and give the omitted vertex label $i$. It follows that
\[
   |O_G|\ge b(D),\qquad |U|\ge p(D)b(D),\qquad |\pi_i(U)|=b(D).
\]
View $U$ as a set of binary words indexed by the unordered pairs of $[n]$. A pair survives deletion of $i$ precisely when neither endpoint is $i$. The $n$ projections therefore cover every coordinate exactly $n-2$ times. Lemma~\ref{lem:proj} gives
\[
 (p(D)b(D))^{n-2}\le |U|^{n-2}
 \le \prod_{i=1}^n |\pi_i(U)| = b(D)^n .
\]
Division proves the bound in terms of $b(D)$. Any parent has at most $n!$ labeled representatives, so $b(D)\le n!$. The asymptotic estimate follows from Stirling's formula.
\end{proof}

\section{Recognizing legitimate decks}
The language \Leg{} consists of the encoded sequences of graphs which form an ordinary vertex deck. It is a decision problem, not the problem of constructing or counting the parents. The counting class $\LWPP$ admits a gap function that is zero on no-instances and has one specified nonzero, polynomial-time computable value, depending only on the input length, on yes-instances~\cite{hhsw}.

\begin{corollary}
\Leg{} belongs to $\LWPP$. In particular, it is not NP-Turing-hard unless NP is low for PP.
\end{corollary}
\begin{proof}
For $n\ge6$, $1\cdot2\cdot3\le n$ and the remaining $n-3$ factors in $n!$ are at most $n$, giving $n!\le n^{n-2}$. Theorem~\ref{thm:main} gives $p(D)\le n^2$. For $3\le n\le5$, the one-card bound $2^{n-1}$ is at most $n^2$. There are at most two parent classes at order two and only one at orders one and zero, if admitted. Thus the single polynomial $q(n)=(n+1)^2$ works at all orders. The assertions now follow directly from Hemaspaandra et al.~\cite[Cor.~4.9--4.10]{hhsw}.
\end{proof}

Membership in $\LWPP$ does not provide a deterministic polynomial-time recognition algorithm. The graph bound counts parents without exhibiting them. It leaves open even whether $p(D)$ is bounded by an absolute constant.

\section{Linearly many pairs with logarithmic neighborhood discrepancy}
\label{geo:section}

For two distinct vertices $u,v$ of a finite simple graph $G$, let
\[
 d_G(u,v)=\bigl|\{w\notin\{u,v\}:\text{$w$ is adjacent to exactly one of $u,v$}\}\bigr|.
\]
Thus $d_G(u,v)$ counts the entries on which their adjacency rows disagree
after the two endpoint positions have been removed. It is not a graph
distance. Write $H_D(G)$ for the graph on $V(G)$ in which $uv$ is an
edge if and only if $d_G(u,v)\le D$; this says nothing about whether
$uv$ is an edge of $G$. The ordinary deck is the multiset of unlabelled
graphs $G-v$, with one occurrence for each vertex. Nonreconstructible
means that some nonisomorphic graph has the same ordinary deck.

\begin{theorem}[A linear matching of near-twins]\label{geo:matching}
Let $G$ be a nonreconstructible graph of order $n\ge 2^{20}$ and put
$D=\lceil320\log_2 n\rceil$. Every vertex cover of $H_D(G)$ has more
than $n/160$ vertices. In particular, every inclusion-maximal matching
of $H_D(G)$ has more than $n/320$ edges.
\end{theorem}

The conclusion concerns distinct \emph{endpoints} in one parent. It does
not say that the sets of vertices distinguishing the matched pairs are
disjoint, or that a surviving mate can be identified on an individual
card. The proof proceeds by contraposition. A small cover leaves a large
set of mutually distinct exterior rows. Sampling exposes enough of their
differences to obtain landmarks. The difficulty is to count all placements
of those landmarks without paying for every permutation of the cover.

\subsection{A landmark criterion requiring no distinguished copy}

A named graph here is a graph on a fixed set of positions. For such a
graph $F$, let $c_X(F)$ be the number of injections of its positions
into $X$ preserving both edges and nonedges. If $|F|<n=|X|$, then
\begin{equation}
 c_X(F)=\frac{1}{n-|F|}\sum_{v\in V(X)}c_{X-v}(F).
 \label{geo:proper}
\end{equation}
Indeed, an injection survives precisely the deletions outside its image.
Thus comparing named pattern counts requires no names on an input card.

Fix an ordered tuple $L$ of $r<n$ distinct vertices of $G$. An exterior
vertex has a binary \emph{word} giving its adjacency to $L$ in order.
Call $L$ resolving if these words are distinct among vertices outside
$L$. No assumption is made about other induced copies of $G[L]$.

\begin{lemma}[Resolving landmarks]\label{geo:landmarks}
Let $L$ be resolving, let $F=G[L]$ be named by the tuple, and let
$N=n-r$. If
\begin{equation}
 \frac{c_G(F)}{|\operatorname{Aut}(G)|}<2^{N/2-1},
 \label{geo:criterion}
\end{equation}
then the ordinary deck determines $G$ up to isomorphism.
\end{lemma}

\begin{proof}
In an arbitrary host $X$ of order $n$, fix one induced injection
$j:F\to X$. Read the exterior words relative to $j$. Discard positions
with a word absent from the target $G$. Every remaining position has a
unique target vertex as its name, though several positions in $X$ may
receive the same name. Define a conflict graph $C_j$ on these positions:
two conflict if their names coincide, or if the names are distinct and
their adjacency in $G$ disagrees with the adjacency of the positions
in $X$. An independent set therefore gives a set $S$ of distinct target
names and exactly one induced injection of $G[L\cup S]$ extending $j$.
Conversely, each such injection gives one independent set in the
conflict graph of its own root placement.

Put $I(C;z)=\sum_{A\text{ independent in }C}z^{|A|}$ and
$P_X(z)=\sum_j I(C_j;z)$, summing over \emph{all} induced placements
of $F$. The bijection just described gives
\begin{equation}
 [z^k]P_X(z)=
 \sum_{\substack{S\subseteq V(G)\setminus L\\|S|=k}}c_X(G[L\cup S]).
 \label{geo:coeff}
\end{equation}
By \eqref{geo:proper}, the coefficients below degree $N$ agree between
hosts with the same deck. At degree $N$ the only possible embeddings are
isomorphisms from the full target. Thus a nonisomorphic host $X$ with
the same deck as $G$ would satisfy
\[
 P_G(z)-P_X(z)=|\operatorname{Aut}(G)|z^N.
\]

Set $b=1+\sqrt2$, $a=2+\sqrt2$. For every graph $C$ on $m$ vertices,
$|I(C;-a)|\le b^m$. For a union of isolated vertices and isolated edges,
use multiplicativity and the factors $1-a=-b$ and $1-2a=-b^2$. Otherwise
split independent sets according to inclusion of a vertex of degree
$d\ge2$. Induction bounds the result by $b^{m-1}+a b^{m-d-1}\le b^m$,
since $b^2+a\le b^3$. This includes the empty graph as the base case.
Consequently $|P_X(-a)|\le c_X(F)b^N$ and likewise for $G$. Since
$c_X(F)=c_G(F)$, the displayed polynomial difference would force
\[
 \frac{c_G(F)}{|\operatorname{Aut}(G)|}
 \ge \frac12(a/b)^N=2^{N/2-1},
\]
contrary to \eqref{geo:criterion}. The tuple $L$ was used only in the
comparison proof; its location was not provided in any card.
\end{proof}

\subsection{Keeping the cover anonymous}

We use one external group-theoretic input. A permutation's support is
the set of points it moves. By Kempe, Pyber and Shalev
\cite[Theorem~A(2)]{KPS2007}, a permutation group on $t\ge2$ points
in which every nonidentity element has support of size at least
$\log_2 t$ has order at most $2^{10t}$. Transitivity is not required.

Let $U\subseteq V(G)$ and $B=V(G)\setminus U$. Partition $U$ by equality
of its full adjacency words on $B$. There is a set $W_0\subseteq B$ of
size at most $|U|$ that defines the same partition. Start with no test
coordinates. Whenever two unequal full words remain together, add a
coordinate on which they differ. Each addition increases the number of
classes. For $U=\varnothing$ take $W_0=\varnothing$.

\begin{lemma}[An anonymous-cover automorphism bound]\label{geo:virtual}
Let $T\supseteq U\cup W_0$, $t=|T|\ge2$, and $B_T=T\setminus U$.
Assume every distinct $b,b'\in B_T$ have at least $\log_2 t$
distinguishers in $G[T]$. Then
\begin{equation}
 \frac{|\operatorname{Aut}(G[T])|}{|\operatorname{Aut}(G)|}
 \le 3^t 2^{10t}.
 \label{geo:quotient}
\end{equation}
\end{lemma}

\begin{proof}
The subgroup $A$ preserving $U$ and $W_0$ separately as sets has index
at most $3^t$: count the possible images of the three-part coloring of
$T$. Its elements permute $B_T$ and permute the classes of $U$, since
equality of words on the setwise preserved $W_0$ detects the full
partition. For each class $C$ create $|C|$ formal positions
$(C,1),\dots,(C,|C|)$. Act on $B_T$ as on actual vertices and on these
new positions by $(C,i)\mapsto(gC,i)$. This is an action on exactly $t$
points. It does not assume that $g$ respects any numbering of the
actual vertices inside $C$.

Its kernel consists precisely of elements fixing $B_T$ pointwise: such
an element fixes $W_0$ pointwise and hence preserves every class. Each
kernel element extends, by the identity off $T$, to an automorphism of
$G$. In fact it fixes the $B_T$ side, and vertices exchanged inside a
class have equal incidences to every vertex of $B$, including those
outside $T$. These extensions are distinct. The kernel therefore has
order at most $|\operatorname{Aut}(G)|$.

A nonidentity element of the image moves some actual $b\in B_T$ to
$b'\ne b$. Any distinguisher in $B_T$ must move. If a vertex of a class
$C$ distinguishes $b,b'$, then all vertices of $C$ do, and $gC\ne C$:
otherwise invariance of the incidences to the class would give the same
bit at $b$ and $b'$. All $|C|$ formal positions of such a class move.
Hence the image element moves at least $d_{G[T]}(b,b')$ points. Apply
the quoted motion bound to the image (the trivial image also satisfies
the order bound), then multiply by the kernel and index estimates.
\end{proof}

\subsection{Sampling the landmarks and closing the count}

\begin{proof}[Proof of Theorem~\ref{geo:matching}]
Suppose $U$ covers $H_D(G)$ and $u=|U|\le n/160$. Put
$\ell=\log_2 n$. Distinct vertices of $B=V(G)\setminus U$ have more
than $D$ distinguishers in $G$, which may include vertices of $U$.
Sample a set $Z$ by including \emph{every} vertex independently with
probability $1/80$. For each fixed pair in $B$, its number $X$ of sampled
distinguishers is binomial with mean $\mu\ge4\ell$. Markov's inequality,
applied to $4^{-X}$, gives
\[
 \Pr(X<\ell)\le 4^\ell\mathbb E[4^{-X}]
 \le \exp\bigl(\ell\ln4-3\mu/4\bigr)
 \le \exp\bigl(-(3-\ln4)\ell\bigr).
\]
A union bound over fewer than $n^2$ pairs is at most $n^{-0.328}<1/26$
for $n\ge2^{20}$, since $(3-\ln4)/\ln2-2>0.328$.
Also $\Pr(|Z|>13n/1000)\le25/26$ by Markov's inequality. Thus there is
one sample for which both bounds hold; independence of these success
events was not used.

Use $W_0$ as above and take $T=U\cup W_0\cup Z$. Then
$t=|T|\le2u+13n/1000\le0.0255n<0.026n$. Every two exterior vertices
are distinguished by $T$, as are all distinct vertices of $B_T$ by at
least $\ell$ vertices of $T$. Thus the exterior words relative to $T$
are distinct and Lemma~\ref{geo:virtual} applies, since $\ell\ge\log_2 t$.
The automorphism bound is immediate also for $t=0,1$.

Temporarily order $T$ and put $F=G[T]$ with these names. Each choice of
its image set supports either no induced injection or precisely
$|\operatorname{Aut}(G[T])|$ of them. By \eqref{geo:quotient},
\[
 \frac{c_G(F)}{|\operatorname{Aut}(G)|}
 \le \binom nt\,3^t 2^{10t}.
\]
With $H_2(x)=-x\log_2 x-(1-x)\log_2(1-x)$ (and the continuous endpoint
values), $\binom nt\le2^{n H_2(t/n)}$ by the binomial theorem. The
base-two logarithm of the right-hand side is at most
\[
 n\bigl(H_2(0.026)+0.026(10+\log_2 3)\bigr)<0.48n
 <(n-t)/2-1.
\]
For the last inequality use $(n-t)/2-1\ge0.487n-1$ in the asserted
range. The landmark criterion \eqref{geo:criterion} now holds with
$N=n-t$, so $G$ is reconstructible. This rules out every such
cover in a counterexample. The endpoints of any inclusion-maximal
matching form a vertex cover: an uncovered edge could be added to the
matching. Dividing the strict cover bound by two proves the claim.
\end{proof}

\section{Editing an arbitrary graph to a reconstructible one}\label{edit:toward}
One edit toggles one unordered pair on a fixed vertex set. Write
$d(G,H)=|E(G)\mathbin{\triangle} E(H)|$, and let $\mathcal R_n$ be the
graphs on $[n]$ uniquely determined, up to isomorphism, by their ordinary
vertex decks. Asilis, Chen, Hansen and Teng~\cite[Section~4, p.~12:15]{acht}
asked about $d(G,\mathcal R_n)=\min_{H\in\mathcal R_n}d(G,H)$. For $n\ge3$, isolating
one vertex gives the bound $n-1$: any graph with an isolated vertex is recovered by
reattaching an isolated vertex to a card of the appropriate edge count.
Their observation about three edits concerns \emph{edge} decks, not vertex
decks.

\begin{theorem}\label{edit:cover}
Let $n\ge 2^{20}$, $D=\lceil320\log_2 n\rceil$ and $k=\lceil2\log_2 n\rceil$.
Every graph $G$ on $[n]$ can be changed to a member of $\mathcal R_n$ in
at most $D+12k$ edits, all incident to at most two vertices. Such an
edited graph can be constructed in polynomial time from $G$.
\end{theorem}

The aim is not just to create twins. We make a twin whose place in an
actual card can be recognized among \emph{all} possible reconstructions.
Call vertices twins if they have identical neighborhoods outside the pair;
their mutual adjacency is unrestricted.

\begin{lemma}\label{edit:twins}
If $n\ge3$ and $2\binom n2<2^{n-2}$, the total number of unordered twin pairs of an
$n$-vertex graph is determined by its ordinary deck.
\end{lemma}
\begin{proof}
Let $s(u,v)$ be the number of vertices outside $\{u,v\}$ adjacent to exactly
one of $u,v$. Form the homogeneous polynomial
\[
 F_G(X,Y)=\sum_{\{u,v\}} X^{s(u,v)}Y^{n-2-s(u,v)} .
\]
Deleting a vertex other than $u,v$ removes one distinguisher or one
nondistinguisher. Thus the sum of $F_J$ over the card occurrences $J$
is $(\partial_X+\partial_Y)F_G$. Any two candidate polynomials differ
by $a(X-Y)^{n-2}$: in coordinates $X-Y,Y$ the differential operator is
just differentiation in the second coordinate. Moreover $a$ is an integer,
as seen from the leading coefficient. If $a\ne0$ the coefficient
$\ell^1$-norm of the difference is at least $2^{n-2}$, whereas the sum
of the two coefficient masses is $2\binom n2$. So $F_G$, in particular
its $Y^{n-2}$ coefficient, is fixed.
\end{proof}

Here is a second ordinary-deck test on a specific card $J$. The total
edge count is the sum of the card edge counts divided by $n-2$; subtracting
$|E(J)|$ gives the missing degree $r$. This determines the entire degree
multiset by ranging over the card occurrences. Write
$Q_G(z)=\sum_{w\in V(G)}z^{\deg_G(w)}$. If the added vertex has
neighborhood $S\subseteq V(J)$, then
\begin{equation}\label{edit:exposure}
 Q_G(z)-Q_J(z)=z^r+(z-1)\sum_{w\in S}z^{\deg_J(w)} .
\end{equation}
Thus both $|S|$ and its neighbor counts in each degree class of $J$ are
prescribed, without matching any vertices across cards.

\begin{proof}[Proof of Theorem~\ref{edit:cover}]
If every pair has $s(u,x)>D$, Theorem~\ref{geo:matching} already gives
$G\in\mathcal R_n$. Otherwise fix a pair with $s(u,x)\le D$. By editing
at most $D$ edges at $x$, make $x$ and $u$ twins. Let $X=G-\{u,x\}$,
write $a_w=\deg_X(w)$, and let $N_0$ be their common neighborhood in
$X$. Retain their mutual bit $e$. Choosing a new common neighborhood
$N$ costs $2|N\mathbin{\triangle}N_0|$ further edits. Call the resulting
graph $P$ and consider $J=P-x$. Put $c=|N|$, $b_w=\mathbf 1_{\{w\in N\}}$.

We shall choose $N$ so that the only possible center of a copy of the
missing twin row on $J$ is $u$. More precisely, a clone of $v$ in $J$
joins every neighbor of $v$ in $J$ and joins $v$ itself according to
a bit $f$. For every $v\in X$ and $f\in\{0,1\}$ we will rule out
its clone by the required degree $c+e$ or by~\eqref{edit:exposure}.

First suppose that the number $s$ of distinct $a_w$ is less than $2k$.
Forbid the at most $3s$ values $a_w,a_w+1,a_w+2$. Choose $c$ in
$[0,n-2]$ within $3s$ of $|N_0|$ such that $c+e$ is not forbidden.
There are at least $3s+1$ available integer choices of $c$ in this
interval even at either end, since $3s<n-1$. Change $N_0$ to a set
of this size by only insertions or only deletions. This costs at most
$6s<12k$ edits. A clone at $v\in X$ has degree $a_v+b_v+f$, one of
the forbidden values.

Now suppose $s\ge2k$. Choose $k$ vertices $W\subset X$ with distinct
$a_w$ separated by at least two: repeatedly take the least remaining
degree and discard it and its successor. Keep $N$ fixed off $W$, allowing
its $k$ bits on $W$ to vary, at cost at most $2k$. For a rival center $v$
and bit $f$, the true exposure sum in~\eqref{edit:exposure} minus the
rival's exposure sum is
\[
 (e-b_v)z^c+(b_v-f)z^{a_v+b_v}
 +\sum_{w\in X\setminus\{v\}}(b_w-A_X(v,w))z^{a_w+b_w},
\]
where $A_X(v,w)$ is the adjacency bit. Fix $f$, $b_v$, and $c$. Changing
$b_w$ from zero to one for $w\in W\setminus\{v\}$ increments the varying
sum by $z^{a_w+1}$ if $A_X(v,w)=0$, and by $z^{a_w}$ otherwise.
These monomials have distinct exponents, so their subset sums are all
different. Among patterns with the fixed triple, at most one can give
zero difference. The total $c$ has at most $k+1$ possible values;
conditioning on it is essential since it depends on the pattern.
Across all rivals, at most $4(n-2)(k+1)<n^2\le2^k$ patterns are ruled
out. Take any of the others.

To verify the certificate against an arbitrary parent, first observe that
$u$ has no twin in $J$. Such a twin $v\in X$ would be interchangeable
with $u$ by an automorphism of $J$, so its clone with $f=e$ would pass
both tests. Every twin pair of $J$ therefore lies in $X$. Its two
vertices see $u$ alike and hence also see $x$ alike in $P$. No old twin
pair was lost, and $u,x$ are twins. Thus $P$ has strictly more twin pairs
than $J$. By Lemma~\ref{edit:twins}, so does every parent of the deck
of $P$. Present any such parent on the same labeled copy of $J$ by
choosing a deletion of that card type. Adding a vertex cannot turn a
surviving nontwin pair into twins, so the new vertex must itself be
a twin of some center in $J$. The two tests exclude every center except
$u$, and the prescribed missing degree fixes $f=e$. Every parent is
therefore isomorphic to $P$.

All pair distances and degree tests above are computable from the
labeled input. In the many-degree case one may try all $2^k<2n^2$
patterns and all their clone tests. This constructs $P$ in polynomial
time, without deciding whether an arbitrary input deck has a parent.
\end{proof}

For smaller orders $n\ge3$ the isolated-vertex bound suffices. Thus the
covering radius of $\mathcal R_n$ is $O(\log n)$; no positive lower
bound is asserted. If the Reconstruction Conjecture holds, this
distance is always zero.

\section{Reconstruction throughout a random edit ball}\label{edit:random}
The preceding theorem chooses favorable edits. We now ask for a graph
that stays reconstructible under \emph{every} edit pattern in a given
ball, even when the pattern is chosen after seeing the graph. Write
$B_b(G)=\{H:d(G,H)\le b\}$.

\begin{theorem}\label{edit:ball}
Let $n\ge2^{20}$, $D=\lceil320\log_2 n\rceil$, and $m=\lceil n/320\rceil$.
Let the edges of $G$ be independent with arbitrary probabilities in
$[\beta,1-\beta]$, where $0<\beta\le1/2$. Set
\[
 q=2\beta(1-\beta),\qquad K=m(n-2m),\qquad
 A=\frac{n!}{(n-2m)!\,2^m m!}.
\]
For every integer $b\ge0$,
\begin{equation}\label{edit:tail}
 \Pr\{B_b(G)\not\subseteq\mathcal R_n\}
 \le \min\{1,\ A\Pr[\operatorname{Bin}(K,q)\le b+mD]\}.
\end{equation}
In particular, if $\beta=\beta(n)$ satisfies $\beta n/\log n\to\infty$,
then for all sufficiently large $n$, uniformly over the allowed edge
probabilities,
\[
 \Pr\{B_{\lfloor c\beta n^2\rfloor}(G)\subseteq\mathcal R_n\}
 \ge1-\exp(-c\beta n^2),\qquad c=10^{-4}.
\]
\end{theorem}
\begin{proof}
Fix a matching $M$ of $m$ unordered vertex pairs (not necessarily graph
edges). For each $\{u,v\}\in M$ and each $w$ outside the $2m$ endpoints,
compare the bits $uw$ and $vw$. Let $L_M$ count the disagreements in
these $K$ comparisons. The comparisons use disjoint pairs of edge
coordinates, so are independent. If the two edge probabilities are
$p,t$, the mismatch probability $p(1-t)+(1-p)t$ is at least $q$,
by minimizing at the four corners of $[\beta,1-\beta]^2$. Thus
$L_M(G)$ stochastically dominates $\operatorname{Bin}(K,q)$.

Each edit changes $L_M$ by at most one. If $H\in B_b(G)$ fails
reconstruction, Theorem~\ref{geo:matching} gives it some matching
$M$ of $m$ pairs, each with at most $D$ distinguishing vertices.
Consequently $L_M(H)\le mD$ and $L_M(G)\le b+mD$. Union-bounding
the fixed-matching probabilities over the $A$ possible $M$ proves
\eqref{edit:tail}. In particular, $M$ need not have been selected
independently of the sampled graph.

For the stated orders, $K\ge0.003n^2$, $q\ge\beta$ and
$\ln A\le2m\ln n$. When $\beta n/\log n\to\infty$, eventually
$mD\le qK/4$ and $\ln A\le qK/16$. If $b\le qK/4$, the
multiplicative Chernoff bound at half the mean gives a failure
probability at most $A\exp(-qK/8)\le\exp(-qK/16)$. The specified
$c$ satisfies both the radius and the exponent bounds.
\end{proof}

For fixed $\beta>0$ this is a quadratic number of arbitrary edits, with
failure exponentially small in $n^2$. Asilis et al.~\cite[Cor.~18--21]{acht}
give an unrestricted edit radius of order $n/\log n$ in this dense
model; their robust asymmetry event has failure $\exp(-\Omega(n))$.
They also allow quadratically many edit locations when the support is
sampled independently of the graph, with edits then unrestricted
within that support~\cite[Lem.~13, Thm.~15]{acht}. Their structural
event controls asymmetry and distinctness of induced subgraphs missing
any fixed number of vertices; their work also yields three-card
reconstruction for the sampled graph. Theorem~\ref{edit:ball} asserts
ordinary full-deck reconstruction, not these stronger properties:
indeed at most $n-2$ edits can make any designated pair twins.


\paragraph{Acknowledgment.} AI tools assisted in developing and checking the argument and in preparing the exposition.

\begin{thebibliography}{9}
\bibitem{acht} J. Asilis, X. Chen, D. Hansen and S.-H. Teng,
Semi-random graphs, robust asymmetry, and reconstruction,
\emph{17th Innovations in Theoretical Computer Science Conference (ITCS 2026)},
LIPIcs \textbf{362} (2026), 12:1--12:21.
\href{https://doi.org/10.4230/LIPIcs.ITCS.2026.12}{doi:10.4230/LIPIcs.ITCS.2026.12}.
\bibitem{KPS2007} J. Kempe, L. Pyber and A. Shalev,
Permutation groups, minimal degrees and quantum computing,
\emph{Groups Geom. Dyn.} \textbf{1} (2007), 553--584.
\href{https://doi.org/10.4171/GGD/24}{doi:10.4171/GGD/24}.
\bibitem{shearer} F. R. K. Chung, R. L. Graham, P. Frankl and J. B. Shearer,
Some intersection theorems for ordered sets and graphs,
\emph{J. Combin. Theory Ser.~A} \textbf{43} (1986), 23--37.
\href{https://doi.org/10.1016/0097-3165(86)90019-1}{doi:10.1016/0097-3165(86)90019-1}.
\bibitem{bt} B. Bollob\'as and A. Thomason,
Projections of bodies and hereditary properties of hypergraphs,
\emph{Bull. London Math. Soc.} \textbf{27} (1995), 417--424.
\href{https://doi.org/10.1112/blms/27.5.417}{doi:10.1112/blms/27.5.417}.
\bibitem{hhsw} E. Hemaspaandra, L. A. Hemaspaandra, H. Spakowski and O. Watanabe,
The robustness of LWPP and WPP, with an application to graph reconstruction,
\emph{43rd International Symposium on Mathematical Foundations of Computer Science (MFCS 2018)}, LIPIcs \textbf{117} (2018), 51:1--51:14.
\href{https://doi.org/10.4230/LIPIcs.MFCS.2018.51}{doi:10.4230/LIPIcs.MFCS.2018.51}.
\bibitem{kh} D. Kratsch and L. A. Hemaspaandra,
On the complexity of graph reconstruction,
\emph{Math. Systems Theory} \textbf{27} (1994), 257--273.
\href{https://doi.org/10.1007/BF01578845}{doi:10.1007/BF01578845}.
\bibitem{wang} J. Wang, Quantitative assembly bounds for vertex-deleted decks: a finite obstruction calculus toward graph reconstruction, preprint (2026).
\href{https://doi.org/10.5281/zenodo.21297717}{doi:10.5281/zenodo.21297717}.
\end{thebibliography}
\end{document}
